% generated from Paper 19/anccft19.tex
\chapter{Algorithmic non-commutative class field theory, XIX: the odd cyclic class that does not exist, and what the exceptional zero sees instead}\label{chap:XIX}
\chaptermark{XIX}
\begin{chapterabstract}
Paper~XIV showed that the finite-dimensional commutative spectral triple of Paper~II has
$HC^1=0$, so no odd cyclic cocycle there can pair to the graph $\cL$-invariant
$-\|h_\alpha\|^2$, and asked whether an infinite-dimensional non-commutative algebra
would help. Three answers, in increasing order of finality. First, the algebra usually
proposed, $\cA_\alpha=C_0(\Yt)\rtimes\Z$, is Morita-trivial: the $\Z$-cover has vertex set
$V\times\Z$ with the deck group acting freely on the second factor, so
$\cA_\alpha\cong\bigoplus_{v\in V}\cK$, with $K_1=0$ and no odd $K$-theory to pair with at
all. Second, the algebra where Bloch theory does live is the $\Z$-equivariant one,
$M_n(C(S^1))$, with $K_1=\Z$; but its natural class, the voltage pencil
$D(t)=(q+1)I-A(t)$, is \emph{not invertible}, because $D(1)$ is the ordinary Laplacian.
On a circle where the pencil is invertible the odd pairing is the winding number of
$\det D$, and that number is the \emph{order} of the exceptional zero --- equal to $2$ for
every voltage class --- whereas $\kappa_0\cL$ is its leading \emph{coefficient},
$\tfrac12\Theta''(0)=-\kappa_0\|h_\alpha\|^2$. Index theory sees orders, not
coefficients. Third, and this is a theorem rather than a diagnosis: an index pairing is
additive in each slot, so if $\alpha\mapsto\tau_\alpha$ and $\alpha\mapsto[u_\alpha]$ are
both linear and $K_1$ has rank one --- as it does for $M_n(C(S^1))$ and for the shadow
algebra alike --- then $\langle\tau_\alpha,[u_\alpha]\rangle$ is a product of two linear
forms in $\alpha$. It cannot be $\kappa_0\|h_\alpha\|^2$, which is positive definite on
$H^1(Y;\R)$, as soon as $b_1(Y)\ge2$. So the impossibility of Paper~XIV is not an artefact
of finite dimensions. Five checks, on the $K_4$ and $K_5$ voltage classes.
\end{chapterabstract}

\section{The proposed algebra is the compacts}\label{XIX:sec:crossed}

Let $Y$ be a finite connected $(q+1)$-regular graph on $n$ vertices, $\alpha$ an integer
$1$-cochain with $[\alpha]\ne0$, and $\Yt$ the associated $\Z$-cover, with deck group
$\Z$. Write $\kappa_0$ for the number of spanning trees of $Y$, $h_\alpha$ for the
harmonic representative of $[\alpha]$, and $\cL(Y_\bullet)=-\|h_\alpha\|^2$
[Ch.~\ref{chap:XIV}, Def.~1.3].

\begin{proposition}\label{XIX:prop:compacts}
The deck group acts freely on the vertex set of $\Yt$ with exactly $n$ orbits, and
\[
\cA_\alpha=C_0(\Yt)\rtimes\Z\;\cong\;\bigoplus_{v\in V}\cK\bigl(\ell^2(\Z)\bigr).
\]
Consequently $K_0(\cA_\alpha)\cong\Z^{\,n}$, $K_1(\cA_\alpha)=0$, and every odd cyclic
cocycle pairs trivially.
\end{proposition}

\begin{proof}
A connected $\Z$-cover of a connected graph has vertex set $V\times\Z$ with the deck group
translating the second coordinate; the action is free and the orbits are the fibres, $n$
of them (check~(A)). Hence $C_0(\Yt)\cong c_0(V)\otimes c_0(\Z)$ equivariantly, with $\Z$
acting only on the second factor, and
$c_0(V)\otimes\bigl(c_0(\Z)\rtimes\Z\bigr)\cong c_0(V)\otimes\cK(\ell^2(\Z))$, the
crossed product of $c_0$ of a free transitive action being the compacts. The
$K$-theory and the vanishing of $HC^{\mathrm{odd}}$ for a direct sum of copies of $\cK$
are standard.
\end{proof}

So the premise that $\cA_\alpha$ ``has nontrivial $K_1$ and $HC^1$'' is false: it is stably
$\C^{\,n}$, and its only cyclic cohomology is the even part, spanned by the $n$ traces.
The non-commutativity of a crossed product by a \emph{free} action is illusory.

\section{The Floquet algebra, and what its pairing sees}\label{XIX:sec:floquet}

The algebra that does carry Bloch theory is the commutant, not the crossed product. After
Floquet transform the $\Z$-equivariant bounded operators on $\ell^2(\Yt)$ are
$M_n\bigl(C(S^1)\bigr)$, whose $K_1$ is $\Z$, generated by the winding of a determinant.
The natural element of it is the voltage pencil $D(t)=(q+1)I-A(t)$ of [Ch.~\ref{chap:IV}, Lem.~1.1],
whose Floquet determinant is $\Theta$.

\begin{proposition}\label{XIX:prop:notinv}
$D$ is not invertible in $M_n(C(S^1))$: $\det D(1)=0$, since $D(1)$ is the ordinary
Laplacian of $Y$. The zero has order exactly two, and
$\tfrac12\Theta''(0)=-\kappa_0\|h_\alpha\|^2$ [Ch.~\ref{chap:IX}, Thm.~1.2]. Hence the class $[D]\in K_1$
does not exist.
\end{proposition}

\begin{corollary}[What the pairing does see]\label{XIX:cor:winding}
Let $\varepsilon$ be small enough that the circle $|t-1|=\varepsilon$ meets no other zero
of $\det D$. There the pencil is invertible and defines a class in $K_1$, and its pairing
with the generator of $HC^1(C^\infty(S^1))$ is the winding number of $\det D$. That
number is the order of the exceptional zero, namely $2$, for \emph{every} voltage class;
it does not depend on $\alpha$, and so cannot be $\kappa_0\|h_\alpha\|^2$.
\end{corollary}

Check~(W) computes this winding number for the three $K_4$ classes and two $K_5$ classes
and finds $2$ in all five, against values of $\kappa_0\|h_\alpha\|^2$ equal to $8$, $16$,
$24$, $75$ and $175$ (Table~\ref{XIX:tab:data}). The contrast is the whole point: an index
pairing is a topological degree, and a degree reads the \emph{order} of vanishing;
$\cL$ is the \emph{coefficient} at the same point.

\begin{table}[ht]
\centering\small
\begin{tabular}{lrrrrr}
\toprule
 & $b_1$ & $\kappa_0$ & $\|h_\alpha\|^2$ & $\kappa_0\|h_\alpha\|^2$ & winding of $\det D$\\
\midrule
$K_4$, $\alpha=(1,0,0)$ & 3 & 16 & $1/2$ & 8 & 2\\
$K_4$, $\alpha=(1,1,1)$ & 3 & 16 & $1$ & 16 & 2\\
$K_4$, $\alpha=(1,2,0)$ & 3 & 16 & $3/2$ & 24 & 2\\
$K_5$, single edge & 6 & 125 & $3/5$ & 75 & 2\\
$K_5$, triangle & 6 & 125 & $7/5$ & 175 & 2\\
\bottomrule
\end{tabular}
\caption{The cleared pencil has $[T^0]=[T^1]=0$ and $[T^2]=-\kappa_0\|h_\alpha\|^2$ in
every case (check~(P)); the winding number is $2$ in every case (check~(W)).}
\label{XIX:tab:data}
\end{table}

\begin{remark}[Where the integrality comes from]\label{XIX:rem:integrality}
The task specification asked why the index should ``absorb'' $\kappa_0=\det G$. It does
not, and no index is involved: by [Ch.~\ref{chap:XIV}, Thm.~7.1], $\|h_\alpha\|^2=a_c^{\mathsf T}G^{-1}a_c$
with $G$ the cycle Gram matrix and $a_c=C^{\mathsf T}\alpha$ integral, so
\[
\kappa_0\|h_\alpha\|^2=a_c^{\mathsf T}\bigl(\det G\cdot G^{-1}\bigr)a_c
=a_c^{\mathsf T}\operatorname{adj}(G)\,a_c\in\Z ,
\]
the adjugate of an integer matrix being an integer matrix. Integrality is Cramer's rule,
not an index theorem (check~(I)).
\end{remark}

\section{The obstruction}\label{XIX:sec:obstruction}

\begin{theorem}[No linear construction, in any algebra with $K_1$ of rank one]
\label{XIX:thm:obstruction}
Let $A$ be a $C^*$-algebra with $\rank_{\Z}K_1(A)=1$. Suppose given group homomorphisms
\[
H^1(Y;\Z)\to HC^1(A),\ \ \alpha\mapsto\tau_\alpha,
\qquad
H^1(Y;\Z)\to K_1(A),\ \ \alpha\mapsto[u_\alpha].
\]
If $b_1(Y)\ge2$ then $\langle\tau_\alpha,[u_\alpha]\rangle\ne\kappa_0\|h_\alpha\|^2$ for
some $\alpha$.
\end{theorem}

\begin{proof}
Torsion classes pair trivially with a $\C$-valued cocycle, so after tensoring with $\Q$
we may write $[u_\alpha]=\ell(\alpha)\,g$ for a fixed generator $g$ of
$K_1(A)\otimes\Q\cong\Q$ and a linear form $\ell$ on $H^1(Y;\Q)$. Additivity of the
pairing in the second slot gives
$\langle\tau_\alpha,[u_\alpha]\rangle=\ell(\alpha)\,\mu(\alpha)$ with
$\mu(\alpha)=\langle\tau_\alpha,g\rangle$, linear by additivity in the first slot. So the
pairing is a product of two linear forms. If $b_1\ge2$ then $\ker\ell\ne0$, so the product
vanishes at some $\alpha$ with $[\alpha]\ne0$; but $\kappa_0\|h_\alpha\|^2>0$ there, since
$h_\alpha=0$ only for $[\alpha]=0$.
\end{proof}

\begin{corollary}\label{XIX:cor:applies}
The theorem applies to $M_n(C(S^1))$, where $K_1\cong\Z$, and to the shadow algebra
$\mathcal O_{B(Y)}$ of [Ch.~\ref{chap:VI}], where $K_1\cong\Z$ by [Ch.~\ref{chap:VI}, Thm.~2.1]. Both candidate
algebras of the task specification are therefore excluded, for every base graph with
$b_1\ge2$ --- in particular for $K_4$ ($b_1=3$) and $K_5$ ($b_1=6$), the graphs on which
the $\cL$-invariant was computed.
\end{corollary}

Check~(Q) verifies the hypothesis that makes the argument bite: $\kappa_0\|h_\alpha\|^2$
is the quadratic form $\operatorname{adj}(G)$, which is positive definite of rank exactly
$b_1$ --- $3$ for $K_4$ and $6$ for $K_5$ --- so it is not a product of two linear forms.

\begin{remark}[What would be needed]\label{XIX:rem:needed}
Theorem~\ref{XIX:thm:obstruction} is not a proof that no odd cyclic class anywhere computes
$\cL$; it rules out the linear constructions in rank-one $K_1$. Two escapes remain, and we
state them as the honest residue of [Ch.~\ref{chap:XIV}, Rem.~9.1(a)]. \emph{Either} one finds an algebra
whose $K_1$ has rank at least $b_1$, with a natural $H^1(Y;\Z)\to K_1$ --- the pairing may
then be a genuine quadratic form; \emph{or} one abandons linearity, in which case the
assignment $\alpha\mapsto\tau_\alpha$ is not a cohomological construction but part of the
input, and the exercise is circular.

The first escape is a real question and it has an obvious candidate, which we record
because the rank matches exactly. The fundamental group of a connected graph is free of
rank $b_1$, and by Pimsner--Voiculescu
\[
K_0\bigl(C^*_r(F_m)\bigr)\cong\Z,\qquad K_1\bigl(C^*_r(F_m)\bigr)\cong\Z^{\,m},
\]
so $C^*_r\bigl(\pi_1(Y)\bigr)=C^*_r(F_{b_1})$ has $\rank K_1=b_1$ --- precisely the
threshold at which Theorem~\ref{XIX:thm:obstruction} stops applying, and the smallest rank at
which a positive definite form on $H^1(Y;\R)$ could be a pairing at all. Whether the
natural map $H^1(Y;\Z)\cong\Z^{b_1}\to K_1(C^*_r(F_{b_1}))\cong\Z^{b_1}$ and a family
$\tau_\alpha$ can be chosen so that the pairing is $\kappa_0\|h_\alpha\|^2$ we do not
know; we note only that $\kappa_0\|h_\alpha\|^2=a_c^{\mathsf T}\operatorname{adj}(G)a_c$ is
already \emph{a} quadratic form on that lattice, so the shape is right and only the
construction is missing.
\end{remark}

\section{Verification}\label{XIX:sec:battery}

\begin{table}[ht]
\centering\small
\begin{tabular}{clc}
\toprule
key & check & mode\\
\midrule
(A) & Proposition~\ref{XIX:prop:compacts}: free deck action, $n$ orbits & exact\\
(P) & the cleared pencil: $[T^0]=[T^1]=0$, $[T^2]=-\kappa_0\|h_\alpha\|^2$ & exact\\
(W) & Corollary~\ref{XIX:cor:winding}: the winding number is $2$ for every $\alpha$ & numeric\\
(I) & Remark~\ref{XIX:rem:integrality}: $\kappa_0\|h_\alpha\|^2=a_c^{\mathsf T}\operatorname{adj}(G)a_c\in\Z$ & exact\\
(Q) & $\operatorname{adj}(G)$ is positive definite of rank $b_1\ge2$ & exact\\
\bottomrule
\end{tabular}
\caption{The verification battery (\texttt{cyclic\_obstruction.py}), $5/5$, about one
second.}
\label{XIX:tab:battery}
\end{table}

\section{Status, scope, corrections}\label{XIX:sec:scope}

\emph{Proved.} Proposition~\ref{XIX:prop:compacts}, Proposition~\ref{XIX:prop:notinv} (given
[Ch.~\ref{chap:IX}, Thm.~1.2]), Corollary~\ref{XIX:cor:winding}, Theorem~\ref{XIX:thm:obstruction} and
Corollary~\ref{XIX:cor:applies}.

\emph{Not proved.} That no odd cyclic class in any algebra computes $\cL$. What is proved
is that the two algebras named, and every algebra with $K_1$ of rank one, fail for a
structural reason, and that the natural class in the Floquet algebra fails for a different
and more interesting one --- it does not exist, and the object obstructing it is the
exceptional zero whose coefficient is the invariant sought.

\begin{remark}[Solved, open, impossible]\label{XIX:rem:soi}
\emph{Solved:} the two candidate algebras of [Ch.~\ref{chap:XIV}, Rem.~9.1(a)] are ruled out, the first
because it is Morita-trivial and the second by Theorem~\ref{XIX:thm:obstruction}.
\emph{Open:} an algebra with $K_1$ of rank $\ge b_1$ and a natural map from $H^1(Y;\Z)$
(Remark~\ref{XIX:rem:needed}); and, unchanged from [Ch.~\ref{chap:XIV}], the question whether $\cL$ has any
index-theoretic description at all, as opposed to the spectral-action and Albanese
descriptions it is already known to have [Ch.~\ref{chap:XIV}, Cor.~7.2].
\emph{Impossible:} an odd pairing linear in $\alpha$ against a rank-one $K_1$, whenever
$b_1\ge2$. Together with [Ch.~\ref{chap:XIV}, Prop.~2.1] this closes the naive routes in both the
finite-dimensional and the infinite-dimensional directions.
\end{remark}

\begin{remark}[Corrections to the task specification]\label{XIX:rem:corrections}
Three. (1) ``$\cA_\alpha=C_0(\Yt)\rtimes\Z$ has nontrivial $K_1$ and $HC^1$'' is false;
the action is free, so the crossed product is a sum of copies of the compacts
(Proposition~\ref{XIX:prop:compacts}). The specification's own hint is the obstruction to
following it. (2) The proposed mechanism --- pair a cocycle built from the integral period
vector $a_c$ with the unit shift $u\in\Z$ and read off $\|h_\alpha\|^2$ from a heat-kernel
limit --- is linear in $\alpha$ in the $K_1$ slot and cannot produce a positive definite
quadratic form (Theorem~\ref{XIX:thm:obstruction}); this is the same objection
[Ch.~\ref{chap:XIV}, Prop.~2.1(iii)] raised against the finite-dimensional proposal, and it survives the
passage to infinite dimensions unchanged. (3) There is nothing for an index to
``absorb'': $\kappa_0\|h_\alpha\|^2=a_c^{\mathsf T}\operatorname{adj}(G)a_c$ is an integer
by Cramer's rule (Remark~\ref{XIX:rem:integrality}), and no index theorem is needed to see it.
\end{remark}

\section*{Provenance of this chapter}

Unrefereed preprint. Every numbered statement is proved above. The computations verify
the exceptional-zero coefficient and the winding number independently on five voltage
classes, and confirm the positive definiteness that makes
Theorem~\ref{XIX:thm:obstruction} bite. The winding number was computed before the theorem was
formulated, in order to find out what the available pairing actually measures.
